消息解码假设 $t\mid q$ 且 $\Delta=q/t$。约束噪声时，对公开乘数 $\Lambda$ 采用固定的系数整数提升。下面的密文相位等式并不声称其分布与新鲜加密相同。

\textbf{- 参考资料：}
\href{https://www.zama.ai/post/tfhe-deep-dive-part-2}{深入理解 TFHE（二）：编码与线性分层运算}~\cite{tfhe-2}

$ $

假设有一个 GLWE 密文 \textsf{ct}：

$\textsf{ct} = \textsf{GLWE}_{S, \sigma}(\Delta M + E) = ( A_0, A_1, \ldots, A_{k-1}, B) \in \mathcal{R}_{\langle n,q \rangle}^{k + 1}$

$ $

\noindent 以及如下新的明文多项式 $\Lambda$：

$\Lambda = \sum\limits_{i=0}^{n-1}(\Lambda_i \cdot X^i) \in \mathcal{R}_{\langle n, q \rangle}$

$ $

\noindent 定义如下密文与明文乘法运算：

$\Lambda \cdot \textsf{ct} = (\Lambda \cdot A_0, \Lambda \cdot A_1, \ldots, \Lambda \cdot A_{k-1}, \Lambda \cdot B)$

$ $

\noindent 当满足变换所需的单位根与可逆性条件时，NTT 乘法使用 $O(n\log n)$ 次域运算（\autoref{sec:ntt}）；例如，满足 $2n\mid q-1$ 的素数 $q$ 支持负循环变换。任意模数并不自动具备这样的 NTT。下面的代数恒等式与所选乘法算法无关：

\begin{tcolorbox}[title={\textbf{\tboxlabel{\ref*{sec:glwe-mult-plain}} GLWE 密文与明文的乘法}}]
$\Lambda \cdot \textsf{GLWE}_{S, \sigma}(\Delta M + E)$

$= \Lambda \cdot (\{A_i^{\langle 1 \rangle}\}_{i=0}^{k-1}, \text{ } B^{\langle 1 \rangle})$

$= (\{\Lambda\cdot A_i^{\langle 1 \rangle}\}_{i=0}^{k-1}, \text{ } \Lambda \cdot B^{\langle 1 \rangle})$

$= \textsf{GLWE}_{S, \sigma}(\Delta (M \cdot \Lambda) + \Lambda\cdot E )$
\end{tcolorbox}

所得相位为 $\Delta M\Lambda+\Lambda E$。只有当 $\|\Lambda E\|_\infty<\Delta/2$ 时，它才解码得到 $M\Lambda\bmod t$。

$ $

%\noindent \textbf{\underline{Proof}}
\begin{myproof}
\begin{enumerate}
\item 定义以下记号： \\
$A_0' = \Lambda \cdot A_0$ \\
$A_1' = \Lambda \cdot A_1$ \\
$\vdots$ \\
$A_{k-1}' = \Lambda \cdot A_{k-1}$ \\
$E' = \Lambda \cdot E$ \\
$B' = \Lambda \cdot B$ \\
\item 可推导得： \\
$B' = \Lambda \cdot B$ \\
$= \Lambda \cdot (\sum\limits_{i=0}^{k-1}{(A_i \cdot S_i)} + \Delta \cdot M + E)$ 
$= \sum\limits_{i=0}^{k-1}{(\Lambda \cdot A_i \cdot S_i)} + \Delta \cdot \Lambda \cdot M + \Lambda \cdot E$ \\ \textcolor{red}{ $\rhd$ 由多项式环的分配律} \\
$= \sum\limits_{i=0}^{k-1}{((\Lambda \cdot A_i) \cdot S_i)} + \Delta \cdot (\Lambda \cdot M) + (\Lambda \cdot E)$  \\
$= \sum\limits_{i=0}^{k-1}{(A_i' \cdot S_i)} + \Delta \cdot (\Lambda \cdot M) + (E')$ \\
\item 由于 $B' = \sum\limits_{i=0}^{k-1}{(A_i' \cdot S_i)} + \Delta \cdot (\Lambda \cdot M) + (E')$，

$(A_0', A_1', \ldots, A_{k-1}'
, B')$ 构成密文 $\textsf{GLWE}_{S, \sigma}(\Delta \cdot \Lambda \cdot M+E^{\prime})$。
\item 因此， \\
$\Lambda \cdot \textsf{GLWE}_{S, \sigma}(\Delta M + E)$ \\
$ = (\Lambda \cdot A_0, \text { } \Lambda \cdot A_1, \ldots, \Lambda \cdot A_{k-1}, \text { } \Lambda \cdot B)$ \\
$ = ( \{A'_{i}\}_{i=0}^{k-1}, \text { } \Lambda \cdot B)$ \\
$= \textsf{GLWE}_{S, \sigma}(\Delta (M \cdot \Lambda) + \Lambda \cdot E)$

%\begin{flushright}
%\qedsymbol{} 
%\end{flushright}
\end{enumerate}
\end{myproof}

使用秘密密钥作减法会精确消去掩码乘积，但舍入只能在解码界内去除 $E'=\Lambda E$。在负循环环中，每个输出系数都是至多 $n$ 个乘积的带符号和，因此
\[
\|\Lambda E\|_\infty\le\|\Lambda\|_1\|E\|_\infty\le n\|\Lambda\|_\infty\|E\|_\infty.
\]
这里 $\|\Lambda\|_1$ 表示所选提升的系数绝对值之和。

重复乘法可能使噪声界增大，直至无法正确解码。解密用于读取结果，并不会刷新已有密文。对于较大的乘数，预先计算的 GLev 表示可以通过下面的 gadget 分解（\autoref{subsec:gadget-decomposition}）给出更好的误差界。

\subsection{利用 gadget 分解抑制噪声}
\label{subsubsec:gadget-decomposition-noise-suppression}

这里使用精确的降幂 gadget $g_i=q/\beta^i$，要求 $q=\beta^l$，并采用 $\{0,\ldots,\beta-1\}$ 中的无符号数位。预先计算的输入是 $\textsf{GLev}^{\beta,l}_{S,\sigma}(\Delta M)$，而非任意单个 GLWE 密文。这样的低噪声 GLev 表示是额外的加密输入，不能假设经过先前运算的密文会自动免费获得它。

在密文与明文乘法 $\Lambda \cdot \textsf{GLWE}_{S, \sigma}(\Delta M)$ 中，噪声 $E$ 增长为 $E' = \Lambda \cdot E$。为限制这种增长，可以对 $\Lambda$ 进行分解（\autoref{subsec:number-decomp}），并使用 $M$ 的 GLev 加密（\autoref{subsec:glev-enc}），具体如下：


$\Lambda = \Lambda_1 \dfrac{q}{\beta^1} + \Lambda_2 \dfrac{q}{\beta^2} + \cdots + \Lambda_l \dfrac{q}{\beta^l} \longrightarrow \textsf{Decomp}^{\beta, l}(\Lambda) = (\Lambda_1, \Lambda_2, \cdots, \Lambda_l)$

$ $

$\textsf{GLev}_{S, \sigma}^{\beta, l}(\Delta M) = \Bigg\{ \textsf{GLWE}_{S, \sigma}\left(\Delta M \dfrac{q}{\beta^1} + E_1\right), \textsf{GLWE}_{S, \sigma}\left(\Delta M \dfrac{q}{\beta^2} + E_2\right), \cdots \textsf{GLWE}_{S, \sigma}\left(\Delta M \dfrac{q}{\beta^l} + E_l\right) \Bigg\}$

$ $


将明文 $M$ 加密为 $\textsf{GLev}_{S, \sigma}^{\beta, l}(\Delta M)$，而非 $\textsf{GLWE}_{S, \sigma}(\Delta M)$；相应地，计算 $\textsf{Decomp}^{\beta, l}(\Lambda) \cdot \textsf{GLev}_{S, \sigma}^{\beta, l}(\Delta M)$，而非 $\Lambda \cdot \textsf{GLWE}_{S, \sigma}(\Delta M)$。注意，两种计算得到的明文结果相同，如下所示：

$\textsf{Decomp}^{\beta, l}(\Lambda) \cdot \textsf{GLev}_{S, \sigma}^{\beta, l}(\Delta M)$

$= (\Lambda_1, \Lambda_2, \cdots, \Lambda_l) \cdot \left (\textsf{GLWE}_{S, \sigma}\left(\dfrac{q}{\beta} \Delta M + E_1\right), \text{ } \textsf{GLWE}_{S, \sigma}\left(\dfrac{q}{\beta^2} \Delta M + E_2\right), \text{ } \cdots, \text{ } \textsf{GLWE}_{S, \sigma}\left(\dfrac{q}{\beta^l} \Delta M + E_l\right) \right )$

$= \Lambda_1\cdot\textsf{GLWE}_{S, \sigma}\left(\dfrac{q}{\beta} \Delta M + E_1\right) +  \Lambda_2\cdot\textsf{GLWE}_{S, \sigma}\left(\dfrac{q}{\beta^2} \Delta M + E_2\right) + \cdots + \Lambda_l\cdot\textsf{GLWE}_{S, \sigma}\left(\dfrac{q}{\beta^l} \Delta M + E_l\right)$

$= \textsf{GLWE}_{S, \sigma}\left(\Lambda_1\cdot\dfrac{q}{\beta}\Delta M + \Lambda_1 E_1\right) +\textsf{GLWE}_{S, \sigma}\left(\Lambda_2\cdot\dfrac{q}{\beta^2}\Delta M + \Lambda_2 E_2\right)+ \cdots + \textsf{GLWE}_{S, \sigma}\left(\Lambda_l\cdot\dfrac{q}{\beta^l}\Delta M + \Lambda_l E_l\right)$

$= \textsf{GLWE}_{S, \sigma}\left(\Lambda_1\cdot\dfrac{q}{\beta} \Delta M + \Lambda_2\cdot\dfrac{q}{\beta^2} \Delta M + \cdots + \Lambda_l\cdot\dfrac{q}{\beta^l} \Delta M+E_{\textit{all}}\right)$

$= \textsf{GLWE}_{S, \sigma}\left(\left(\Lambda_1\cdot\dfrac{q}{\beta} + \Lambda_2\cdot\dfrac{q}{\beta^2} + \cdots + \Lambda_l\cdot\dfrac{q}{\beta^l}\right)\cdot \Delta M  + E_{\textit{all}} \right)$ \textcolor{red}{ $\rhd$ 其中 $E_{\textit{all}} = \sum\limits_{i=1}^l \Lambda_iE_i$}

$= \textsf{GLWE}_{S, \sigma}\left(\Lambda \cdot \Delta M + E_{\textit{all}}\right)$ \textcolor{red}{ $\rhd$ 解密结果为 $\Lambda\cdot M$}

$ $

精确分解时，最终误差为 $E_{\textit{all}}=\sum_{i=1}^l\Lambda_iE_i$。若 $\|E_i\|_\infty\le B_{\mathrm{g}}$，无符号数位给出
\[
\|E_{\textit{all}}\|_\infty\le\sum_{i=1}^l\|\Lambda_i\|_1\|E_i\|_\infty\le nl(\beta-1)B_{\mathrm{g}}.
\]
若要使用 $\beta/2$ 的数位界，必须明确改用平衡分解。这可能改善大乘数对应的最坏情形误差界，但不能保证每个输入的实际误差都更小。只有各自满足误差界时，两种计算才会解码得到相同消息。对于截断分解 $\Lambda\equiv\sum_i\Lambda_i g_i+\rho\pmod q$，最终相位的误差变为 $\sum_i\Lambda_iE_i-\rho\Delta M$；还必须约束这一残差。图~\ref{fig:decomp2} 展示了预期的误差界改善。

\begin{figure}[h!]
    \centering
  \includegraphics[width=0.8\linewidth]{figures/decomp2.pdf}
  \caption{利用 gadget 分解降低密文与明文乘法中的噪声。}
  \label{fig:decomp2}
\end{figure}

\subsubsection{讨论}
\label{subsubsec:glwe-mult-plain-discuss}


分解技术还具有重要用途：在 GLWE 密钥切换（\autoref{sec:glwe-key-switching}）中，我们将介绍如何将分解后的掩码多项式 $\textsf{Decomp}^{\beta,l}(A_i)$
与预先计算的密钥切换密钥
$\textsf{KSK}_i=\textsf{GLev}^{\beta,l}_{S',\sigma}(S_i)$,
结合起来实现密钥切换。因此，尽管每次操作输出的都是普通 GLWE 密文，gadget 分解仍可在多次密钥切换调用中反复使用。

对于反复重新初始化普通密文噪声 $E$ 的技术，我们将在 \autoref{subsec:tfhe-noise-bootstrapping} 介绍 TFHE 的噪声自举。
